%==============================================================================
%  Scientific poster — 36 x 48 in (portrait) — tikzposter
%  "Finding Gravitational Lenses from Super-Resolved Euclid Images"
%  Aleksandr Belotserkovtsev
%
%  Layout: LEFT COLUMN = vertical pipeline infographic (top->bottom);
%          RIGHT = one text column (abstract -> results).
%
%  Compile: pdflatex poster -> biber poster -> pdflatex poster -> pdflatex poster
%  (references come from references.bib via biblatex). Notes:
%   - figure placeholders are fixed-size boxes (replace \figbox with
%     \includegraphics once images are available)
%   - references are a short manual list (swap for biblatex + references.bib
%     if you prefer automatic citations)
%==============================================================================
\documentclass[25pt, a0paper, portrait,
               margin=0mm, innermargin=15mm,
               blockverticalspace=8mm, colspace=14mm, subcolspace=8mm]{tikzposter}

% --- Custom paper size: 36 x 48 inches (overrides a0paper above) -------------
\geometry{paperwidth=36in, paperheight=48in}

\usepackage[T1]{fontenc}
\usepackage[utf8]{inputenc}
\usepackage{amsmath, amssymb}
\usepackage{graphicx}
\usepackage{lmodern}
\usetikzlibrary{positioning, arrows.meta, fit, backgrounds, calc}

% --- Bibliography (biblatex + biber) -----------------------------------------
% Compile: pdflatex poster -> biber poster -> pdflatex poster -> pdflatex poster
\usepackage[backend=biber, style=authoryear, sorting=nyt, maxbibnames=2,
            giveninits=true, doi=false, isbn=false, url=false]{biblatex}
\addbibresource{references.bib}
\renewcommand*{\bibfont}{\footnotesize}

% --- Theme & colours ---------------------------------------------------------
\usetheme{Default}
\usecolorstyle[colorPalette=BlueGrayOrange]{Denmark}

% --- Figure placeholder: \figbox{height}{label} ------------------------------
% Replace with \includegraphics[width=\linewidth]{file} when assets are ready.
\newcommand{\figbox}[2]{%
  \begin{center}%
  \setlength{\fboxsep}{0pt}%
  \fcolorbox{gray!55}{gray!12}{%
    \parbox[c][#1][c]{0.96\linewidth}{\centering\textcolor{gray!65}{\Large #2}}}%
  \end{center}}

% --- Pipeline panel: \panelimg{width}{file}{fallback label} ------------------
% Uses the rendered image if present, else a sized gray placeholder. Render
% panels with:  python3 scripts/poster_render.py <your.fits> --out fig/poster
\newcommand{\panelimg}[3]{%
  \IfFileExists{#2}{\includegraphics[width=#1]{#2}}{%
    \setlength{\fboxsep}{0pt}%
    \fcolorbox{gray!55}{black!88}{%
      \parbox[c][#1][c]{#1}{\centering\textcolor{gray!60}{\small #3}}}}}

% --- Pipeline panel widths (set relative to the pipeline column linewidth) ----
\newlength{\imgs}     % side of a single square panel (rows 1-2)
\newlength{\lanegap}  % horizontal gap between the left and right lane
\newlength{\fullw}    % rows 3-4 span both lanes: \fullw = 2\imgs + \lanegap

%==============================================================================
%  Title
%==============================================================================
\title{\parbox{0.95\linewidth}{\centering Finding Gravitational Lenses from
        Super-Resolved Euclid Images}}
\author{Aleksandr Belotserkovtsev}
\institute{Center for Astrophysics \,$|$\, Harvard \& Smithsonian \\[4pt]
           \normalsize Mentors: Asst.\ Prof.\ L.\ Connor \quad Dr.\ Barone \quad
           AGEL Collaboration}

\begin{document}
% Title bar spans the full paper width (margin=0mm makes \paperwidth edge-to-edge).
\maketitle[width=\paperwidth, roundedcorners=0, linewidth=0pt, innersep=20pt]

%==============================================================================
%  Single row, three columns:
%    col 1 (0.46) = vertical PIPELINE infographic  (the "huge" figure)
%    col 2 (0.27) = text  |  col 3 (0.27) = text
%==============================================================================
\begin{columns}

  %========================= COLUMN 1 — PIPELINE ==============================
  \column{0.50}
  \block{The Pipeline}{%
  \centering
  % ---- swap these paths for your own thoughtfully-chosen renders ----
  % Render with: python3 scripts/poster_render.py <your.fits> --out fig/poster
  \def\figdir{fig/poster}
  \vspace{6mm}
  % Pipeline geometry: two vertical guide lines bound the figure. Rows 1-2 hold
  % two equal squares (\imgs) left/right-aligned to the guides; rows 3-4 hold one
  % image spanning both guides (\fullw = 2\imgs + \lanegap).
  \setlength{\imgs}{0.30\linewidth}
  \setlength{\lanegap}{0.05\linewidth}
  \setlength{\fullw}{2\imgs}\addtolength{\fullw}{\lanegap}
  \begin{tikzpicture}[
      >={Stealth[length=11mm, width=8mm]},
      panel/.style={inner sep=0pt, outer sep=0pt},
      % Side captions: left-lane labels sit LEFT of the image, the star-cluster
      % and ePSF labels sit RIGHT of it. Big, minimal text.
      ltitle/.style={font=\LARGE\bfseries, text=black, align=right,
                     anchor=east, inner sep=0pt},
      rtitle/.style={font=\LARGE\bfseries, text=black, align=left,
                     anchor=west, inner sep=0pt},
      flow/.style={->, line width=2.6pt, gray!70},
      opl/.style={font=\small\bfseries, text=white, fill=black!80,
                  inner sep=5pt, rounded corners=4pt},
      node distance=10mm]

    % Each multi-panel stage is a SINGLE composite image built by
    % scripts/poster_grid.py (no hand-stacking of cells in LaTeX). The left
    % lane is west-aligned to x=0; the right lane is east-aligned to the right
    % guide; rows 3-4 span both guides.

    % ===== ROW 1: ingredients (left) + star-cluster map (right) =====
    \node[panel, anchor=north west] (ing) at (0,0)
        {\panelimg{\imgs}{\figdir/ingredients_grid.png}{ingredients grid}};
    \node[ltitle] at ([xshift=-6mm]ing.west) {Ingredients};
    \node[panel, anchor=north west] (cmap) at ([xshift=\lanegap]ing.north east)
        {\panelimg{\imgs}{\figdir/cluster_map.png}{star-cluster map}};
    \node[rtitle] at ([xshift=6mm]cmap.east) {Star\\ clusters};

    % ===== ROW 2: clean field (left) + per-band PSF grid (right) =====
    \node[panel, anchor=north west] (clean) at ([yshift=-22mm]ing.south west)
        {\panelimg{\imgs}{\figdir/clean_field.png}{clean synthetic field}};
    \node[ltitle] at ([xshift=-6mm]clean.west)
        {Synthetic\\ field\\[3pt]{\normalsize\color{gray!75} ground truth}};
    \node[panel, anchor=north west] (psf) at ([yshift=-22mm]cmap.south west)
        {\panelimg{\imgs}{\figdir/psf_grid.png}{per-band PSFs}};
    \node[rtitle] at ([xshift=6mm]psf.east) {Per-band\\ ePSF};

    % ===== ROW 3: mock-Euclid dirty grid spans BOTH guides =====
    \node[panel, anchor=north west] (dirty) at ([yshift=-28mm]clean.south west)
        {\panelimg{\fullw}{\figdir/dirty_grid.png}{mock Euclid (4 bands)}};
    \node[ltitle] at ([xshift=-6mm]dirty.west)
        {Mock\\ Euclid\\[3pt]{\normalsize\color{gray!75} 4 noisy bands}};

    % ===== ROW 4: super-resolved image spans BOTH guides =====
    \node[panel, anchor=north west] (sr) at ([yshift=-24mm]dirty.south west)
        {\panelimg{\fullw}{\figdir/sr_field.png}{super-resolved}};
    \node[ltitle] at ([xshift=-6mm]sr.west)
        {Super-\\ resolved\\[3pt]{\normalsize\color{gray!75} 0.05$''$/pix}};

    % ===== FLOW ARROWS =====
    \draw[flow] (ing.south)  -- (clean.north);                      % ingredients -> clean
    \draw[flow] (cmap.south) -- (psf.north);                        % clusters -> PSF grid
    \draw[flow] (clean.south) -- (dirty.north);                     % clean -> dirty
    \draw[flow] (psf.south)   -- (dirty.north);                     % PSFs  -> dirty
    \node[opl] at ([yshift=8mm]dirty.north) {$\circledast$ \,$+$\, noise};
    \draw[flow] (dirty.south) -- (sr.north) node[opl, midway] {POLISH\texttt{++} CNN};
  \end{tikzpicture}
  }

  % References + Acknowledgements live under the pipeline (back-matter, uses the
  % space below the figure so the single text column carries only the body).
  \block{References}{
    \nocite{*}% list every entry in references.bib, not only those cited inline
    \printbibliography[heading=none]
  }

  \block{Acknowledgements}{
    \small Conducted with mentorship from the Center for Astrophysics
    $|$ Harvard \& Smithsonian and the AGEL collaboration.
    \\[4pt]\textbf{Contact:} abelotserkovtsev@college.harvard.edu
  }

  %========================= COLUMN 2 — TEXT (single) ========================
  \column{0.50}

  \block{Abstract}{
    Strong gravitational lensing yields systems that are valuable because they
    allow for studies of distant magnified galaxies that would otherwise be too
    faint to see, the mass distribution of the foreground galaxy or cluster,
    dark matter, and even the geometry of the Universe. In particular, lenses
    with \emph{small Einstein radii} (the angular distance between the lens and
    the source) are especially valuable, since they constrain the dark-matter
    halo profile. However, when this radius is too small, the lens cannot be
    visually separated from the source in low spatial resolution, wide-area
    surveys, which causes underrepresentation of lenses with small Einstein
    radii in studies. Recent advances in super-resolution techniques offer a way
    to recover this missing population by improving the identification of strong
    gravitational lenses in \textbf{Euclid} surveys, the first wide-area
    space-based optical survey. We build on \textbf{POLISH}
    \parencite{Connor2022} and \textbf{POLISH++} \parencite{Wu2026} architectures,
    originally developed for radio astronomy: they are deep convolutional neural
    networks with a global residual connection. We adapt POLISH to increase the
    resolution of Euclid images to the level of higher-quality \textbf{Hubble}
    images, training on synthetic Euclid-like images built from simulated
    galaxies paired with a high-resolution ground truth, and benchmarked against
    Hubble imaging. Our
    preliminary results show that, applied to real Euclid data, the model both
    super-resolves \emph{and} denoises, revealing structure that is blurred or
    buried in the originals --- a promising first step toward recovering the
    missing small-radius lenses.
  }

  \block{Background}{
    A massive foreground galaxy bends the light of a background source into
    arcs or multiple images; the \emph{Einstein radius} $\theta_E$ sets their
    angular separation. When $\theta_E$ falls below the survey resolution the
    arcs blend into the lens light --- those systems are the ones that best
    constrain dark-matter halos, and the ones wide surveys miss.
    \vspace{0.5em}
    \figbox{8cm}{[ gravitational lens ]}
  }

  \block{Model Architecture}{
    \textbf{WDSR-A} (POLISH/POLISH++): a deep \textbf{CNN} with wide-activation
    residual blocks, weight-normalised convolutions, a \textbf{global residual
    skip}, and $\times2$ pixel-shuffle upsampling. Input $4$ bands
    (VIS\,+\,Y/J/H), output a deconvolved VIS sky. No batch-norm; trained with an
    asinh-stretched L1 loss.
  }

  \block{Significance \& Next Steps}{
    \begin{itemize}
      \item Applies cutting-edge computer-vision super-resolution to optical
            astronomy --- not for precise photometry, but to flag lens
            \emph{candidates} for follow-up.
      \item Targets the underrepresented \textbf{small-$\theta_E$} population,
            highly valuable for cosmology and galaxy evolution.
      \item Deliverables: open-source simulation code, trained weights, an
            end-to-end super-resolution pipeline, and a paper.
    \end{itemize}
  }

  \block{Preliminary Results: Real Euclid Data}{
    Applied to \textbf{real Euclid cutouts}, the model both
    \textbf{super-resolves and denoises}, recovering structure that is blurred
    or buried in the originals.
    \vspace{0.5em}
    \begin{tikzfigure}
      \begin{minipage}{0.47\linewidth}
        \figbox{9cm}{[ Euclid input ]}
      \end{minipage}\hfill
      \begin{minipage}{0.47\linewidth}
        \figbox{9cm}{[ SR output ]}
      \end{minipage}
    \end{tikzfigure}
    \vspace{0.3em}
    Benchmarked against Hubble imaging of the same targets.
  }

\end{columns}

\end{document}
